\subsection{Pure cubes and squarefree triples cannot contribute an extremal boundary point}

\begin{proposition}[P-F29: the complementary squarefree-support bound]
\label{prop:cubicsquarefreesupport}
Let $n\ge2$, and let a cubic $F\in\mathbb C[x_1,\ldots,x_n]$
have only monomials of the forms $x_i^3$ and $x_ix_jx_k$
with $i,j,k$ pairwise distinct. Let $B$ be its support
exponent matrix, of rank $r$ over $\mathbb Q$, and assume
$\mathbf1\in\operatorname{im}_{\mathbb Q}B$. Let $l$ be the
number of variables whose cubes occur in the support.
If a diagonal matrix $A$ of finite two-power order fixes
$F$ exactly, then
\[
 \operatorname{ord}(\det A)\le
 \begin{cases}
  3^{(r-l)/2},&r<n\text{ or }l>0,\\
  3^{n/2-1},&r=n\text{ and }l=0.
 \end{cases}
 \tag{P-F29a}
\]
In both cases this is strictly less than $2^{n-1}$.

Consequently, every nonzero competing fixed quotient point
in the discriminant boundary of P-F26 has, in every
eigenbasis of every finite two-power exact-form lift with
determinant order $2^{n-1}$, both a monomial $x_i^2x_j$
with $i\ne j$ and a monomial in three pairwise distinct
variables.
\end{proposition}

\begin{proof}
As in P-F27, the eigenvalue character
$\psi:\mathbb Z^n\to\mathbb C^*$ of $A$ annihilates
the lattice generated by the support columns, and
$\psi(\mathbf1)=\det A$. Its image is a two-group.
The $l$ cube columns $3e_i$ are linearly independent.
Include all of them in a selection of $r$ support columns
forming a basis of the rational column span, and write
$C$ for the resulting $n\times r$ matrix. There is a
unique rational vector $c$ with
\[
 Cc=\mathbf1.
 \tag{P-F29b}
\]

Choose a nonzero $r\times r$ minor $C_I$ by selecting
rows. Those rows must contain all $l$ cube-variable
indices: otherwise a cube column would be zero in the
minor. After reordering rows and columns, expansion in
the cube columns gives
\[
 \det C_I=\pm3^l\det E,
 \tag{P-F29c}
\]
where $E$ is an $(r-l)\times(r-l)$ matrix with entries
zero or one and at most three ones in each column.
For $r=l$ the empty determinant is interpreted as one.
Hadamard's inequality gives
\[
 0<|\det E|\le3^{(r-l)/2}.
 \tag{P-F29d}
\]
For clarity, this use of Hadamard is elementary: the
absolute determinant is the product of the successive
orthogonalized column lengths, each at most the original
Euclidean column length, which here is at most $\sqrt3$.

Let $q$ be the least positive integer such that $qc$ is
integral. Cramer's rule in $C_Ic=\mathbf1_I$ shows that
$q$ divides $|\det C_I|$. Also
$q\mathbf1=C(qc)$ belongs to the support lattice, so
the order of $\psi(\mathbf1)$ divides $q$.
Because that order is a power of two, the odd factor
$3^l$ in (P-F29c) does not contribute, and
\[
 \operatorname{ord}(\det A)
       \le 2^{v_2(\det E)}
       \le |\det E|\le3^{(r-l)/2}.
 \tag{P-F29e}
\]
If $r<n$ or $l>0$, then $r-l\le n-1$.
Since $\sqrt3<2$ and $n-1>0$, this is strictly less
than $2^{n-1}$, proving the first case. This reasoning
allows a rank-deficient support matrix and an arbitrary
two-power order for the matrix $A$ itself.

In the remaining case $r=n$ and $l=0$, the selected
matrix $C$ is square and every column has exactly three
ones. Thus
$\mathbf1^tC=3\mathbf1^t$, so reduction modulo three
makes its rows dependent and $3\mid\det C$.
The two-part of a nonzero integer divisible by three
is at most one third of its absolute value. Therefore
\[
 \operatorname{ord}(\det A)
    \le 2^{v_2(\det C)}
    \le\frac{|\det C|}{3}
    \le3^{n/2-1}<2^{n-1}.
 \tag{P-F29f}
\]
The last strict inequality follows immediately for
$n=2$, and for $n>2$ from
$3^{(n-2)/2}<2^{n-2}<2^{n-1}$.
This proves (P-F29a).

For the quotient consequence, P-F26 supplies the finite
two-power lifts and P-F27 supplies the rational-span
hypothesis from semistability in every coordinate basis.
Absence of a three-distinct-variable monomial would put
the support in P-F27's graph class and force a smooth
chain, contrary to discriminant vanishing. Absence of
an $x_i^2x_j$ monomial would put it in the support class
of this proposition and contradict the primitive
determinant order. Hence both types must occur.
\end{proof}

\paragraph{Remaining support class.}
This is a complementary exclusion, not a classification
of arbitrary singular semistable cubics. The remaining
representatives must mix square-linear monomials with
three-distinct-variable monomials. By the connected-graph
extension of P-F27, their square-linear graph must also
be disconnected. The interaction of triple monomials
with these different graph components is the remaining
exponent-lattice problem. No historical priority is
claimed for the elementary determinant estimate.
